\section*{Complexity Analysis}

This section analyzes the cost of the operations introduced in the previous section.

\subsection*{Cost of the primitive operations}

The basic subroutine is evaluation of a represented expression.

\begin{lemma}
    Evaluating the expression represented by $G$ at a fixed assignment takes $O(|G|)$ time.
\end{lemma}
\begin{proof}
  The evaluation algorithm processes the nodes of $G$ in topological order.
  Each node is visited once, and each visit performs only constant-time work beyond combining the already computed values of its children.
\end{proof}

\subsection*{Cost of the tests from the substitution-map representation}

\begin{lemma}
  Let $p$ be a represented polynomial with DAG $G$.
  \begin{enumerate}
    \item Testing whether $p \equiv 0$ by evaluating it at $c$ sampled points takes $O(|G|)$ time.
    \item Testing whether $p$ is linear in a chosen variable $x_1$ takes $O(|G|)$ time.
    \item If $p$ is linear in $x_1$, recovering the represented coefficients $a$ and $b$ in
      \[
        p(x_1,\dots,x_n) = a x_1 + b
      \]
      by the identities $b = p(0,x_2,\dots,x_n)$ and $a = p(1,x_2,\dots,x_n) - p(0,x_2,\dots,x_n)$ takes $O(|G|)$ time.
  \end{enumerate}
\end{lemma}
\begin{proof}
  For the zero test, we perform $c$ evaluations of $G$, one at each sampled point thus the total cost is $O(|G| \cdot c) = O(|G|)$.

  For the linearity test, we perform a zero test for a polynomial with a constant number of additional nodes, thus the total cost is $O(|G|)$.

  Coefficient recovery involves a constant number of copies and evaluations of $G$, thus the total cost is $O(|G|)$.
\end{proof}

\subsection*{One elimination step and overall growth}

Let $N(Q,\Sigma)$ denote the total number of DAG nodes appearing in the equations in $Q$ and the substitution map $\Sigma$.
We assume that arithmetic in the base field and memoized node lookup both take constant time.

\begin{lemma}
  One elimination step starting from $(Q_i,\Sigma_i)$ has cost $O(N_i)$, where $N_i = N(Q_i,\Sigma_i)$.
\end{lemma}

\begin{proof}
Let $(p,x)$ be the chosen pair. By the previous lemma, testing whether $p$ is linear in $x$ and recovering the coefficients needed to solve for $x$ cost $O(|p|)$. Forming $p(0)$ and $p(1)$ and building the represented substitution
\[
  x \mapsto -\frac{b}{a}
\]
add only $O(|p|)$ work, since they duplicate or evaluate $p$ a constant number of times. Removing the active equations that have become zero costs $O(N_i)$, because each zero test is linear in the size of its DAG and the total size of the active equations is at most $N_i$. Since $|p| \leq N_i$, the whole step costs $O(N_i)$.
\end{proof}

Moreover, the substitution is stored symbolically rather than expanded through all active equations, so one successful elimination step adds only a constant number of represented expressions derived from $p$. Hence there is a constant $K$ such that
\[
  N_{i+1} \leq N_i + K|p| \leq (1+K)N_i.
\]

\begin{theorem}
  The total running time of the elimination phase is
  \[
    O\!\left(\sum_{i=0}^{d-1} N_i\right) = O(N_0 c_0^d).
  \]
\end{theorem}
\begin{proof}
  First note that for $c_0 = 1 + K$ we have $N_{i+1} \leq c_0 N_i$ for every successful elimination step; iterating gives $N_d \leq c_0^d N_0$. 
  Since each step costs $O(N_i)$, the total runtime is $O(\sum_{i=0}^{d-1} N_i) = O(N_0 c_0^d)$.
\end{proof}

In summary, all basic operations are linear in the size of the current represented state.
Each successful elimination step adds only a constant amount of new represented structure, because substitutions are kept symbolic instead of being expanded throughout the system.
As a result, the represented state grows by at most a constant factor per step, and the total runtime is exponential in the elimination depth.
